\documentclass[11pt,a4paper]{article}
\usepackage[utf8]{inputenc}
\usepackage[T1]{fontenc}
\usepackage[portuguese]{babel}
\usepackage{geometry}
\geometry{margin=2.5cm}
\usepackage{amsmath,amssymb}
\usepackage{booktabs}
\usepackage{listings}
\usepackage{xcolor}
\usepackage{enumitem}
\usepackage{hyperref}
\hypersetup{colorlinks=true,linkcolor=blue!50!black,urlcolor=blue!50!black}

\lstset{
  basicstyle=\ttfamily\small,
  breaklines=true,
  frame=single,
  framesep=4pt,
  rulecolor=\color{black!30},
  numbers=none,
  showstringspaces=false,
}

\title{Vigilância de Partições Retangulares\\
\large Projeto Prático --- MAD CC3003 2025/26}
\author{Grupo 06\\
Cauã Pinheiro Souza (202302423) \and Matheus Gonçalves Guerra (202205324)}
\date{Maio 2026}

\begin{document}
\maketitle

\section{Introdução}

O problema consiste em determinar uma colocação mínima de guardas
nos vértices de uma partição retangular $\Pi$, garantindo que todos
os retângulos de $\Pi$ (ou um subconjunto $\Pi' \subset \Pi$ dado)
ficam vigiados. Um guarda colocado num vértice $p$ vigia todos os
retângulos incidentes em $p$.

\paragraph{Notação.} Seguimos a notação do problema dos
``Postos de Vigia'' (slides MAD2526\_clp e clp\_swi\_2026):
\begin{itemize}[noitemsep]
  \item $R$ --- retângulos a cobrir.
  \item $P$ --- vértices candidatos a posto de guarda.
  \item $G[r]$ --- vértices que veem o retângulo $r$.
  \item $V[p]$ --- retângulos que o vértice $p$ vê (inversa de $G$).
  \item $S \subseteq P$ --- conjunto de guardas escolhidos (solução).
\end{itemize}

\paragraph{Pré-processamento partilhado por todas as abordagens.}
\begin{enumerate}[noitemsep,topsep=2pt]
  \item Leitura do ficheiro produzido pelo gerador \texttt{rectParts.c}
    (formato: na linha~1, o número de instâncias; cada instância
    começa pelo número de retângulos $nr$, seguida de $nr$ linhas
    com \texttt{id $\;$ nv $\;$ x1 y1 \dots xnv ynv}).
  \item Identificação do retângulo exterior como aquele com a
    \emph{bounding box} máxima.
  \item Descarte dos vértices da borda externa do candidato a guarda,
    pela observação da alínea~(b) do exercício~7 da fp3:
    para qualquer solução ótima, é sempre possível trocar um guarda
    da borda por um vizinho interno equivalente.
\end{enumerate}


\section{Q1 --- Estratégias \emph{greedy}}

\subsection{Métodos}

Avaliámos duas estratégias gulosas, ambas com origem direta no
material da cadeira.

\paragraph{(1) \emph{first-fail}.} Em cada iteração escolher o
retângulo $r$ não coberto com menos candidatos disponíveis (menor
$|G[r] \setminus S|$); entre os candidatos desse retângulo,
escolher o que cobre mais retângulos ainda não cobertos. A relação
com a heurística \emph{first-fail} usada na procura de soluções de
CSPs é discutida explicitamente na alínea~(a) do exercício~7
da fp3 e nos slides MAD2526\_clp: atacar primeiro a variável mais
restringida (com menor domínio) é onde a procura é mais provável
falhar e portanto onde mais se corta o espaço de procura. No nosso
problema, o ``domínio'' de cobertura de cada retângulo é o conjunto
$G[r]$ dos vértices que o veem.

\paragraph{(2) \emph{max-cobertura}.} Em cada iteração escolher o
vértice $p$ que maximiza $|V[p] \cap \text{não\_cobertos}|$. É a
estratégia clássica para o problema de \emph{minimum set cover} e
segue o mesmo padrão das estratégias gulosas apresentadas nos
slides (Kruskal: escolher o ramo de menor peso; \emph{knapsack}
fraccionário: escolher pelo maior rácio valor/peso) --- escolha
baseada num critério local de maior benefício por passo.

\paragraph{Complexidade.} Ambas correm em $O(|R|\cdot|P|)$, polinomial
em $n=|R|+|P|$.

\paragraph{Pseudo-código \emph{first-fail}.}
\begin{lstlisting}
S, nao_cobertos <- {}, R
enquanto nao_cobertos != {}:
    r* <- argmin_{r em nao_cobertos} |G[r] \ S|
    p* <- argmax_{p em G[r*] \ S} |V[p] cap nao_cobertos|
    S <- S cup {p*};  nao_cobertos <- nao_cobertos \ V[p*]
\end{lstlisting}

\subsection{Avaliação experimental}

Para avaliar a qualidade das heurísticas, comparámo-las contra o
ótimo obtido por Programação Inteira (modelo dos slides \mbox{MAD2526\_clp},
\textit{Postos de Vigia}, resolvido com Google OR-Tools / SCIP).

Geraram-se baterias de instâncias com o gerador \texttt{rectParts.c}
(distribuição geométrica, $\tau = 0{,}75$) para
$nr \in \{10, 20, 40, 80, 160\}$, com 20 instâncias por tamanho.

\begin{table}[h]
\centering
\small
\setlength{\tabcolsep}{4pt}
\begin{tabular}{rrrrrrrrrr}
\toprule
$nr$ & ot.\ méd. & ff méd. & ff \%ót. & ff exc.\ \% & mc méd. & mc \%ót. & mc exc.\ \% & $t_{\text{PI}}$ (ms) \\
\midrule
 10 & 3{,}80 & 4{,}00 & 80{,}0 \% & 6{,}25\,\% & 4{,}15 & 65{,}0 \% & 10{,}00\,\% & 4{,}6 \\
 20 & 7{,}20 & 7{,}85 & 50{,}0 \% & 9{,}29\,\% & 8{,}55 & 0{,}0 \% & 18{,}93\,\% & 5{,}5 \\
 40 & 14{,}40 & 16{,}20 & 0{,}0 \% & 12{,}65\,\% & 17{,}15 & 5{,}0 \% & 19{,}27\,\% & 13{,}5 \\
 80 & 27{,}60 & 31{,}50 & 0{,}0 \% & 14{,}14\,\% & 33{,}60 & 0{,}0 \% & 21{,}78\,\% & 24{,}4 \\
 160 & 55{,}25 & 63{,}05 & 0{,}0 \% & 14{,}15\,\% & 68{,}80 & 0{,}0 \% & 24{,}54\,\% & 102{,}0 \\
\bottomrule
\end{tabular}
\caption{Q1: número médio de guardas (ót.\ = ótimo;
ff = \emph{first-fail}; mc = \emph{max-cobertura}). ``\%ót.'' = percentagem de
instâncias em que a heurística atinge o ótimo. ``exc.\,\%'' = excesso médio
relativo ao ótimo.}
\label{tab:q1}
\end{table}

\paragraph{Validação na fp3 ex.~7.} Ambas as heurísticas atingem o ótimo
conhecido ($|S| = 4$) na instância canónica, em parte porque essa
instância contém retângulos com $|G[r]|=1$, que forçam a escolha.

\subsection{Discussão}

\paragraph{\emph{first-fail} domina \emph{max-cobertura}.} Em todos os
tamanhos a estratégia \emph{first-fail} produz menos guardas. A diferença
estabiliza em cerca de 10 pontos percentuais de excesso médio nas
instâncias maiores ($\approx 14\,\%$ vs $\approx 24\,\%$ para $nr=160$).
Embora a \emph{max-cobertura} seja a heurística clássica para
\emph{set cover}, com garantia de aproximação $H(|R|) = O(\log |R|)$,
ela falha aqui porque ignora a estrutura geométrica do problema: ao
maximizar a cobertura local pode deixar retângulos com poucos candidatos
sem opções boas.

\paragraph{Estrutura das instâncias.} Numa análise das instâncias
geradas, verificou-se que praticamente nunca há retângulos com
$|G[r]|=1$ (i.e., guardas forçados), mas existe uma cauda de
retângulos com $|G[r]| \in \{2,3\}$. São estes que o \emph{first-fail}
ataca primeiro, capturando a maior parte do ganho heurístico.

\paragraph{Tempo.} O tempo das heurísticas é desprezável (poucos ms);
o solver exato continua rápido (102~ms em média para $nr=160$).
Para os tamanhos estudados não há vantagem em usar heurísticas em
detrimento do ótimo; a sua utilidade reside na compreensão da
estrutura do problema e como termo de comparação para variantes
mais complexas (Q4).


\section{Q2 --- Programação Inteira e Programação por Restrições}

\subsection{Modelo Matemático}

Para modelar o problema de vigilância de partições retangulares como
um problema de otimização combinatória, partilhamos a mesma base
teórica entre a Programação Inteira (PI) e a Programação por
Restrições (CSP). Recorrendo ao pré-processamento estabelecido na Q1
(onde a borda externa é descartada), reutilizamos o conjunto de
vértices candidatos válidos $P$ e o conjunto de retângulos a cobrir
$R$. O modelo define-se da seguinte forma:

\begin{itemize}[topsep=2pt,itemsep=2pt]
  \item \textbf{Variáveis de decisão:} $x_p \in \{0,1\}$ para cada
    $p \in P$, onde $x_p = 1$ indica a colocação de um guarda no
    vértice $p$.
  \item \textbf{Função objetivo:} minimizar o número total de guardas
    colocados na partição,
    \[ \min \sum_{p \in P} x_p . \]
  \item \textbf{Restrições de cobertura:} para garantir que todos os
    retângulos são vigiados, cada retângulo $r \in R$ tem de possuir
    pelo menos um guarda num dos seus vértices incidentes $G[r]$,
    \[ \sum_{p \in G[r]} x_p \ge 1 \qquad \forall\, r \in R . \]
\end{itemize}

\subsection{Implementação MAC (\emph{Maintaining Arc Consistency}) e \emph{Branch \& Bound}}

Na alínea~(b) implementámos um algoritmo de pesquisa com retrocesso
(\emph{backtracking}) que não recorre a motores externos de CSP. Para
garantir eficiência computacional, o algoritmo foi reforçado com duas
técnicas cruciais:

\begin{itemize}[topsep=2pt,itemsep=2pt]
  \item \textbf{Propagação de restrições (MAC):} em cada nó da árvore
    de pesquisa aplicamos consistência de arco. Se, após algumas
    atribuições, um retângulo ficar sem guardas atribuídos e com
    apenas um vértice livre restante no seu conjunto $G[r]$, o
    algoritmo força imediatamente a colocação de um guarda nesse
    vértice ($x_p = 1$). Se um retângulo ficar sem guardas e sem
    vértices livres, o ramo falha imediatamente.
  \item \textbf{\emph{Branch and bound} (\emph{pruning}):} o algoritmo
    mantém em memória o custo da melhor solução viável encontrada
    ($S^*$). Qualquer ramo da árvore de pesquisa que acumule um número
    de guardas igual ou superior a $|S^*|$ é imediatamente podado,
    reduzindo drasticamente o espaço de procura.
\end{itemize}

\subsection{Soluções com motores externos: CP-SAT e Prolog (clpfd)}

Para as alíneas seguintes, o mesmo modelo foi traduzido para
ferramentas \emph{state-of-the-art}:

\begin{itemize}[topsep=2pt,itemsep=2pt]
  \item \textbf{Google OR-Tools (CP-SAT):} o modelo foi mapeado
    utilizando a interface \texttt{cp\_model}, instanciando vetores de
    variáveis booleanas. As restrições foram aplicadas diretamente
    iterando sobre as listas de índices dos vértices incidentes a cada
    retângulo, tirando partido da inferência lógica do solver SAT
    subjacente.
  \item \textbf{SWI-Prolog (clpfd):} a mesma lógica foi declarada em
    Prolog usando domínios finitos (\texttt{ins 0..1}). A cobertura de
    cada face foi garantida com a restrição
    \texttt{sum(VarsDaFace, \#>=, 1)}, seguida de um predicado de
    otimização \texttt{labeling([min(MinGuardas)], Guardas)}.
\end{itemize}


\section{Q3 --- Programação Dinâmica}

Para que a aplicação desta técnica valha a pena, devemos avaliar se é
possível repartir o problema em subproblemas independentes e se isso
compensa face às restantes soluções possíveis.

\subsection{A dificuldade da subestrutura em grafos 2D arbitrários}

Ao modelarmos a partição retangular num grafo (onde os retângulos são
os elementos a cobrir e os vértices são os pontos de decisão),
obtemos uma estrutura muito densa e com diversos ciclos.

A decisão de colocar um guarda num vértice $v_i$ resolve a cobertura
dos retângulos incidentes, mas altera o ``estado'' necessário para
tomar decisões nos vértices vizinhos. Como o mapa é bidimensional,
estas decisões propagam-se em várias direções, fechando ciclos e
destruindo a independência local exigida pela Programação Dinâmica
clássica (como a aplicada em árvores ou sequências lineares). O estado
necessário para memorizar um subproblema não se resume a um único nó,
mas sim a toda a fronteira de intersecção com o resto do mapa.

\subsection{Programação Dinâmica de Perfil (\emph{Profile DP})}

A única forma exata de aplicar Programação Dinâmica a este problema de
grelha bidimensional seria através de uma técnica conhecida como
Programação Dinâmica de Perfil (\emph{Profile DP} ou varrimento de
linha). O algoritmo funcionaria, em teoria, da seguinte forma:

\begin{enumerate}[topsep=2pt,itemsep=0pt]
  \item Uma linha imaginária varre o mapa ortogonal da esquerda para a
    direita (ou de cima para baixo).
  \item A linha interseta, a qualquer momento, um conjunto de
    retângulos, formando uma ``fronteira'' ou ``perfil'' de largura
    $W$.
  \item O estado da PD teria de memorizar todas as combinações
    possíveis de cobertura para os retângulos dessa fronteira.
  \item Para cada nova coluna avançada, o algoritmo testaria todas as
    permutações válidas de guardas e faria a transição de estados.
\end{enumerate}

\subsection{Complexidade}

A complexidade temporal e espacial da Programação Dinâmica de Perfil é
governada pelo tamanho máximo da fronteira $W$. O número de estados a
memorizar em cada passo cresce exponencialmente, na ordem de
$\mathcal{O}(2^W)$. No nosso caso de estudo (por exemplo, uma
instância com 40 retângulos e 82 vértices interligados num espaço
denso), a largura do perfil $W$ atinge valores que originam uma
explosão combinatória abrupta na tabela de memorização. Conclui-se
que, embora matematicamente possível, esta abordagem é
computacionalmente impraticável e significativamente inferior às
técnicas baseadas em restrições utilizadas na Q2.


\section{Q4 --- Extensões}

\subsection{Q4a --- Guardas com cores}

\paragraph{Enunciado e modelo.}
Cada guarda ativo recebe uma cor. Requisito: se dois guardas $p$
e $q$ veem ambos o mesmo retângulo $r$ (i.e.~$p,q \in G[r]$),
as suas cores diferem. Em linguagem CSP/CLP é uma restrição
\texttt{all\_distinct} aplicada às cores dos guardas ativos em cada
$G[r]$ --- a mesma primitiva usada nos slides MAD2526\_clp (SEND+MORE,
Sudoku) e em clp\_swi\_2026 (\textit{quatroEstilos}).

Equivalentemente, define-se o \emph{grafo de conflito} $H = (S, E)$
com $E = \{\{p,q\}: \exists r,\ p,q \in S \cap G[r]\}$; o número
mínimo de cores é o número cromático $\chi(H)$.

\paragraph{Duas leituras da pergunta.}
A pergunta do enunciado --- ``para o número mínimo de guardas,
qual seria o número mínimo de cores?'' --- admite duas interpretações,
porque o número mínimo de guardas pode ser atingido por várias
soluções diferentes, cada uma com um grafo de conflito distinto:
\begin{description}[topsep=2pt,itemsep=0pt]
  \item[Leitura A (sequencial).] Fixar uma solução ótima de guardas
    $S$ e calcular $\chi(H_S)$ --- o número mínimo de cores para
    \emph{essa} colocação.
  \item[Leitura B (lexicográfica).] Entre \emph{todas} as colocações
    com o número mínimo de guardas $G^*$, escolher a que minimiza o
    número de cores.
\end{description}
Implementámos ambas. A Leitura~B fornece uma resposta mais forte à
pergunta; a Leitura~A é a interpretação mais direta e útil quando
a solução de guardas já está fixada por outras razões.

\paragraph{Modelo de Leitura B (lexicográfico).} Variáveis
booleanas $y_{p,k}$ = ``guarda $p$ ativo com cor $k$'' e
$\text{usa}_k$ = ``cor $k$ usada por algum guarda''.
\begin{align*}
\min \quad & \textstyle\sum_k \text{usa}_k \\
\text{s.a.} \quad & \textstyle\sum_{p \in G[r]} x_p \ge 1, \quad r \in R \\
& \textstyle\sum_p x_p = G^* \\
& \textstyle\sum_k y_{p,k} = x_p, \quad p \in P \\
& \textstyle\sum_{p \in G[r]} y_{p,k} \le 1, \quad r \in R,\ k \\
& y_{p,k} \le \text{usa}_k \\
& \text{usa}_k \ge \text{usa}_{k+1} \quad \text{(quebra de simetria)} \\
& x_p, y_{p,k}, \text{usa}_k \in \{0,1\}
\end{align*}
A restrição $\sum_{p \in G[r]} y_{p,k} \le 1$ é exatamente a tradução
para variáveis booleanas da restrição \texttt{all\_distinct} sobre
as cores dos guardas ativos em $G[r]$. Resolvido com CP-SAT
(Google OR-Tools).

\paragraph{Resultados experimentais (30 instâncias por tamanho).}
\begin{table}[h]
\centering
\small
\begin{tabular}{rrrrrrrr}
\toprule
$nr$ & guardas méd. & cores$_A$ méd. & $A$ máx. & cores$_B$ méd. & $B$ máx. & $B$ prov.\ ót. \\
\midrule
 10 & 3{,}9 & 1{,}90 & 3 & 1{,}13 & 2 & 100\,\% \\
 20 & 7{,}4 & 2{,}13 & 3 & 1{,}37 & 2 & 100\,\% \\
 40 & 14{,}2 & 2{,}10 & 4 & 1{,}60 & 3 & 100\,\% \\
 80 & 27{,}7 & 2{,}23 & 4 & 1{,}63 & 2 & 100\,\% \\
\bottomrule
\end{tabular}
\caption{Q4a: número de cores nas duas leituras.}
\label{tab:q4a}
\end{table}

\paragraph{Discussão.}
\begin{itemize}[topsep=2pt,itemsep=0pt]
  \item O número de cores mantém-se baixo em todas as instâncias
    estudadas ($\le 4$ na Leitura~A, $\le 3$ na Leitura~B). A
    explicação geométrica é que o número de cores é limitado
    inferiormente pelo tamanho da maior \emph{clique} de guardas que
    se veem todos no mesmo retângulo, e cada retângulo tem um número
    limitado de vértices candidatos.

  \item A Leitura~B dá sistematicamente menos cores do que a Leitura~A
    (média $\approx 1{,}5$ contra $\approx 2{,}1$). Em muitas instâncias
    basta uma cor na Leitura~B, indicando que existe uma colocação
    ótima de guardas em que nenhuns dois guardas veem o mesmo retângulo.

  \item Resposta à pergunta do enunciado: para o número mínimo de
    guardas, o número mínimo de cores é tipicamente $1$ ou $2$ e
    raramente $3$. Depende criticamente de qual solução ótima de
    guardas se escolhe --- é o ponto subtil que distingue as duas
    leituras.
\end{itemize}

\subsection{Q4b --- Guardas com maior alcance}

\paragraph{Enunciado e interpretação.}
``Permitir que um guarda num vértice possa guardar não só os
retângulos incidentes mas também vizinhos a uma distância (no grafo)
não superior a um valor $D$ dado como parâmetro.''

A decisão central é a interpretação de ``o grafo''. Adotámos a
seguinte: \emph{o grafo é o grafo dos vértices da subdivisão planar}
--- nós são os vértices $(x,y)$ da partição; arestas são os
segmentos entre vértices consecutivos ao longo das linhas da grelha.
Justificação:
\begin{enumerate}[noitemsep,topsep=2pt]
  \item A frase do enunciado diz ``um guarda \emph{num vértice} \dots
    distância (no grafo)'' --- a distância é naturalmente medida a
    partir do vértice onde o guarda está, no grafo onde esse
    vértice vive.
  \item Para $D=0$ (``vizinhos a distância 0'' $=$ o próprio vértice)
    recupera-se exatamente a definição base ``vê as peças incidentes
    nesse vértice'' (fp3 ex.~7). Verificámos isto computacionalmente.
  \item O ficheiro \texttt{types.h} do gerador define adjacência entre
    vértices (campo \texttt{nxt[2]}: ``próximo vértice na mesma linha
    XX e YY''). É a estrutura de grafo presente no código.
\end{enumerate}

\paragraph{Modelo modificado.} Idêntico ao modelo base, com
$G[r]$ substituído por
\[
  G_D[r] \;=\; \{\, p \in P : \exists u,\ \text{dist}_{\text{grafo}}(p,u) \le D
  \ \text{e}\ u \in G[r]\,\}.
\]

\paragraph{Algoritmo.} O grafo não tem pesos, pelo que a distância
mede-se em número de arestas. Calculamos $V_D[p]$ por
\emph{BFS limitado a profundidade $D$} (slide DAA2526\_grafos5),
caso particular de Dijkstra com pesos unitários
(slide DAA2526\_grafos3). O solver de cobertura é o mesmo da Q2a.

\paragraph{Resultados experimentais (30 instâncias por tamanho).}
\begin{table}[h]
\centering
\small
\begin{tabular}{rrrrrr}
\toprule
$nr$ & $D=0$ & $D=1$ & $D=2$ & $D=3$ & $D=4$ \\
\midrule
 10 & 3{,}9 & 2{,}1 & 1{,}8 & 1{,}1 & 1{,}0 \\
 20 & 7{,}4 & 4{,}3 & 3{,}0 & 2{,}0 & 1{,}8 \\
 40 & 14{,}2 & 7{,}9 & 5{,}1 & 3{,}3 & 2{,}5 \\
 80 & 27{,}7 & 15{,}2 & 9{,}5 & 6{,}3 & 4{,}5 \\
\bottomrule
\end{tabular}
\caption{Q4b: número médio de guardas em função do alcance $D$.}
\label{tab:q4b}
\end{table}

\paragraph{Discussão.}
\begin{itemize}[topsep=2pt,itemsep=0pt]
  \item O alcance reduz substancialmente o número de guardas
    necessários. Em $nr=80$, passar de $D=0$ para $D=1$ corta o
    número de guardas quase para metade ($27{,}7 \to 15{,}2$).
    Cada incremento adicional continua a reduzir, com retornos
    decrescentes; para $D$ suficientemente grande, um único
    guarda chega.

  \item Verificação: para $D=0$ o número médio de guardas coincide
    exatamente com o valor reportado na Q1 (Tabela~\ref{tab:q1}),
    confirmando que a extensão generaliza corretamente o problema
    base.

  \item Os tempos de resolução continuam baixos
    --- o gargalo computacional é o solver de cobertura, não o BFS.
\end{itemize}


\section{Conclusão}

Resumo das contribuições:
\begin{itemize}[topsep=2pt,itemsep=0pt]
  \item Na Q1, a estratégia \emph{first-fail} sugerida pela
    alínea~(a) do exercício~7 da fp3 traduz-se diretamente num
    algoritmo guloso eficaz, consistentemente superior à
    estratégia clássica de \emph{max-cobertura}, com excesso médio
    estabilizado em $\approx 14\,\%$ relativamente ao ótimo nas
    instâncias maiores.
  \item Na Q2, a avaliação cruzada entre a implementação manual (MAC)
    e os motores externos (CP-SAT e Prolog) validou a exatidão do
    modelo matemático, retornando invariavelmente os mesmos custos
    ótimos. Destaca-se a observação empírica da simetria de soluções
    ótimas, em que diferentes solvers atingem o mesmo número mínimo de
    guardas, mas com alocações geográficas distintas devido às
    heurísticas de procura próprias de cada motor.
  \item Na Q3, a análise teórica e prática da Programação Dinâmica
    demonstrou que a densidade do grafo bidimensional planar e a
    imprevisibilidade da largura da fronteira ($W$) geram uma explosão
    combinatória exponencial ($\mathcal{O}(2^W)$). Conclui-se que esta
    abordagem, embora matematicamente possível, é computacionalmente
    impraticável e significativamente inferior às técnicas baseadas em
    restrições utilizadas na Q2.
  \item Na Q4a, distinguir as duas leituras da pergunta do enunciado
    revela um aspeto subtil: o número mínimo de cores depende de qual
    solução ótima de guardas se considera. Em ambas as leituras o
    número de cores é pequeno ($\le 4$) e na Leitura~B é tipicamente
    $1$ ou $2$.
  \item Na Q4b, com o grafo dos vértices como leitura natural de
    ``o grafo'', o número de guardas decresce rapidamente com $D$
    e a generalização do modelo é direta.
\end{itemize}

\appendix
\section{Estrutura da entrega}

\begin{lstlisting}
GRUPO_06/
  relatorio.pdf                 (este documento)
  README.txt                    (procedimentos de execucao)
  programas/
    rectParts.c, types.h, macros.h   (gerador da professora)
    Q1/
      instance.py    -- leitura e pre-processamento
      greedy.py      -- estrategias first-fail e max-cobertura
      exato.py       -- solver de PI (OR-Tools, referencia otima)
      benchmark.py   -- estudo experimental da Q1
    Q2/
      instance_q2.py -- leitura e pre-processamento
      ex2B_MAC_AC3.py-- MAC (AC-3) com branch & bound
      ex2C_ORTools.py-- solver CP-SAT (OR-Tools)
      ex2C_prolog.pl -- solver clpfd (SWI-Prolog)
      ex2A.md        -- modelo matematico
    Q3/
      instance_q3.py -- leitura e pre-processamento
      ex3_DP.py      -- programacao dinamica (bitmask)
      ex3.md         -- avaliacao da aplicabilidade da PD
    Q4/
      instance_q4.py -- leitura com ordem e construcao do grafo
      alcance.py     -- Q4b (set cover com alcance D)
      cores.py       -- Q4a (Leituras A e B)
      bench_alcance.py
      bench_cores.py
  casos_teste/
    inst_10.txt, inst_20.txt, inst_40.txt, inst_80.txt
\end{lstlisting}

\end{document}
